\input{estilo}
\begin{document}
\begin{figura}\begin{tikzpicture}[pieza/.style={caja, minimum width=24mm, minimum height=9mm}]
  % dos nodos iguales: worker, chunk server y el disco local
  \foreach \n/\x in {1/0, 2/5.2} {
    \draw[marco] (\x,0) rectangle (\x+3.2,-6.2);
    \node[above=1pt] at (\x+1.6,0) {nodo \n};
    \node[pieza] (w\n) at (\x+1.6,-0.9) {worker};
    \node[pieza] (cs\n) at (\x+1.6,-2.7) {chunk server};
    \pic[scale=1.7] at (\x+1.6,-5.85) {base};
  }
  % en el nodo 1, el chunk local entra como input del worker
  \node[archivo, minimum width=6.5mm, minimum height=8mm] (ch) at (1.15,-4.8) {};
  \node[anchor=west, font=\small] at (ch.east) {64 MB};
  \draw[->] (ch.north) -- (cs1.south -| ch.north);
  \draw[->] (cs1) -- node[etiq, right] {input} (w1);
  % el shuffle va de worker a worker, sin pasar por GFS
  \draw[->] (w1) -- node[above, etiq] {shuffle} (w2);
\end{tikzpicture}\end{figura}
\end{document}
